% back-nozzle.tex

\newpage
\section{Flow through a conical nozzle}
\label{back-nozzle-sec}
%
Good quality experimental data for wall pressure distribution in a conical
nozzle with a circular-arc throat profile and a 15$^o$ divergent section
is available in Ref.\,\cite{back_etal_1965a}.
In the original experiment, the flow of air through the facility 
was allowed to reach steady state and static pressures were measured 
at a large number of points along the nozzle wall.

% In contrast, the present simulation is transient with just the transonic 
% plus supersonic parts of the flow field reaching steady state.
% Fixed 2014

Figure\,\ref{back-geometry-fig} shows the outline of the simulated flow
domain which is set up to approximate the largest subsonic area ratio
used in the experiment.
A short subsonic section upstream of the throat is included, along with
the conical supersonic expansion where the pressure measurements were made.
Note that the geometric calculation of the tangent arcs is done within the
input script.
This makes use of Python, beyond just being an input format, and allows the
specification to be fully parametric.
Although the parametric description makes the initial setup of the script 
a bit more complex than absolutely necessary, 
it does make the running of the simulation for other radii of curvature 
very simple.

\begin{figure}[htbp]
\begin{center}
\includegraphics[width=\textwidth]{../2D/back-nozzle/back-layout.pdf}
\end{center}
\caption{Schematic diagram of the full flow domain for the duct and conical nozzle.}
\label{back-geometry-fig}
\end{figure}

% The relatively long upstream part of the simulated tube provides the gas 
% through an unsteady expansion from zero speed and pressure of 500\,kPa 
% (state 4) up to a small Mach number (state 3).
% These state labels refer to the those for 
% the hypothetical shock tube problem in which 
% state 1 is the initial low-pressure condition, 
% state 4 is the initial high-pressure condition,
% state 2 is the post-shock condition of the low-pressure gas,
% and state 3 is the expanded high-pressure gas condition.

% Along the $C_+$ characteristic connecting states 4 and 3, 
% the Riemann invariant can be written as 
% $$
%     J_+ = u + \frac{2 a}{\gamma - 1} ~~~,
% $$
% and so the relation between states 4 and 3 can be written as
% $$
%     \frac{a_4}{a_3} = 1 + \frac{\gamma - 1}{2} M_3 ~~~.
% $$
% Because the unsteady expansion is isentropic and $a = \sqrt{\gamma R T}$,
% the pressure ratio can be written as
% $$
%     \frac{p_4}{p_3} = \left[ 1 + \frac{\gamma - 1}{2} M_3 
%                       \right]^{2 \gamma / (\gamma - 1)}
% $$
% which gives a specific pressure ratio of $p_4/p_3 = 1.2102$.
% Since the experiment used a steady expansion from a large reservoir
% at stagnation conditions to the equivalent of our state 3, 
% the corresponding stagnation pressure for state 3 is computed from
% $$
%     \frac{p_{03}}{p_3} = \left[ 1 + \frac{\gamma - 1}{2} M_3^2
%                                \right]^{\gamma / (\gamma - 1)} 
% $$
% which gives $p_{03} = 418.7$\,kPa in the current simulation.

Figure\,\ref{back-mesh-fig} shows the mesh, coloured by Mach number 
(once the flow has reached steady state).
Assuming that flow in the subsonic and transonic regions of the nozzle is
steady, the expected Mach number is $M_3 = 0.13812$ for an area ratio of 
$A_3 / A_* = 4.2381$.
This is seen to be consistent with the Mach number colouring in the figure
and is a good test of the \verb!SubsonicInBC! that is applied at the 
upstream boundary.\index{boundary conditions!SubsonicInBC!example of use}

\begin{figure}[htbp]
\begin{center}
\includegraphics[width=0.8\textwidth]{../2D/back-nozzle/back-mach-field-with-mesh.png}
\end{center}
\caption{Mesh generated for the axisymmetric nozzle simulation, coloured with Mach number.}
\label{back-mesh-fig}
\end{figure}

Figure\,\ref{back-pressure-contours-fig} shows the pressure distribution
throughout the flow domain at $t = 4.0$\,ms, once the flow has settled.
Note that the inflow boundary has the flow stagnation properties specified as its
flow condition but that this condition does not appear in any part of the simulation domain.

\begin{figure}[htbp]
\begin{center}
\includegraphics[width=0.8\textwidth]{../2D/back-nozzle/back-p-field.png}
\end{center}
\caption{Pressure contours within the flow domain at 4.0\,ms.}
\label{back-pressure-contours-fig}
\end{figure}

The flow in the nozzle is largely transient as the stagnation conditions drive gas into the domain
but the overall flow becomes steady, as indicated by the histories
shown in Fig.\,\ref{back-transient-fig}.
Because there is little damping to the gas dynamics, small scale oscillations evident in the pressure
history take some to damp out as weak waves bounce around in the subsonic region long after the bulk flow
has approached steady state.
Figure\,\ref{back-profile-fig} shows that the simulation matches the
experimental data closely.

\begin{figure}[htbp]
\begin{center}
\begin{tabular}{cc}
(a) & (b) \\
\includegraphics[width=0.5\textwidth]{../2D/back-nozzle/back_history_M.pdf} &
\includegraphics[width=0.5\textwidth]{../2D/back-nozzle/back_history_p.pdf}
\end{tabular}
\end{center}
\caption{Development of the flow at a ``history point'' near
         the centre of the exit plane:
         (a) Mach number;
         (b) static pressure.}
\label{back-transient-fig}
\end{figure}

\begin{figure}[htbp]
\begin{center}
\begin{tabular}{cc}
(a) & (b) \\
\includegraphics[width=0.5\textwidth]{../2D/back-nozzle/back_profile_whole.pdf} &
\includegraphics[width=0.5\textwidth]{../2D/back-nozzle/back_profile_supersonic.pdf}
\end{tabular}
\end{center}
\caption{Normalised pressure distribution along the nozzle wall:
         (a) full length of flow domain; 
	 (b) just the supersonic part of the nozzle.}
\label{back-profile-fig}
\end{figure}

\newpage
\subsection{Input script (.py)}
\topbar
\lstinputlisting[language={}]{../2D/back-nozzle/back.py}
\bottombar


\subsection{Shell scripts}
\label{back-sh-files}
\topbar
\lstinputlisting[language={}]{../2D/back-nozzle/back_run.sh}
\bottombar

\noindent
\topbar
\lstinputlisting[language={}]{../2D/back-nozzle/back_profile.sh}
\bottombar

\noindent
\topbar
\lstinputlisting[language={}]{../2D/back-nozzle/back_history.sh}\index{history location!extracting the data}
\bottombar

\newpage
\subsection{Notes}
\begin{itemize}
\item The simulation reaches a final time of 4\,ms in 5410 steps and,
      on an AMD Phenom II X4 840 system, this takes 128 seconds run time.
\item The pressure is normalised with respect to the stagnation pressure
      using the following AWK script.\\
      \topbarshort
      \lstinputlisting[language={}]{../2D/back-nozzle/normalize.awk}
      \bottombarshort\\
\item The history data for all of the history cells in a particular block are
      written to the one file.
      A particular cell can be extracted as shown by the following AWK script.\\
      \topbarshort
      \lstinputlisting[language={}]{../2D/back-nozzle/extract-history.awk}
      \bottombarshort\\
\end{itemize}
